% \chapter{Literature review}
\chapter{LITERATURE REVIEW}
\label{ch:lit review}

This chapter reviews prior work on automated exam evaluation, 
the use of large language models (LLMs) in education, and cognitive 
assessment frameworks. The thesis does not aim to build a new grading 
system but to test how reliable LLM-generated results are when compared 
with teacher evaluations. The review focuses on three areas: 
(1) applications of LLMs in educational assessment, (2) their reasoning 
and text-understanding abilities, and (3) approaches for validating AI 
results against human judgments. The chapter closes with Bloom’s 
taxonomy as a framework for aligning both AI and teacher assessments.

\section{The Integration of LLMs in Educational Assessment}

% With the rapid development of LLMs, researchers have begun to 
% explore their role in automating feedback and assessment, 
% especially for open-ended and subjective tasks. 
% \cite{Holmes2022} argue that LLMs can enhance large-scale 
% education by providing consistent, immediate feedback while 
% simultaneously raising concerns regarding cultural bias and 
% fairness in underrepresented languages. This issue is 
% particularly relevant for Thai, where available training data 
% is limited compared to high-resource languages such as English 
% or Chinese.
\subsection{State of the art and practice in AI in education}

As highlighted by \cite{Holmes2022}, the accelerating 
advancement of Large Language Models (LLMs) has intensified 
interest in their potential to automate educational 
processes—particularly feedback and assessment. Early studies 
demonstrated that LLMs can analyze written responses and 
provide formative feedback with a level of fluency and 
contextual understanding previously unattainable by 
rule-based systems. The integration of artificial 
intelligence (AI) into education, or Artificial Intelligence 
in Education (AIED), 
can be categorized into three major domains: student-focused, 
teacher-focused, and institution-focused applications. Each 
category serves distinct purposes but shares the goal of 
improving learning outcomes and administrative efficiency.

Student-focused applications include intelligent tutoring 
systems, adaptive learning platforms, and personalized 
feedback mechanisms that respond dynamically to each 
learner’s needs. Teacher-focused tools encompass automated 
grading systems, plagiarism detection, and analytics 
dashboards that help educators identify student 
misconceptions or engagement levels. Institution-focused 
systems, meanwhile, support large-scale educational 
management, such as admissions decision-making, e-proctoring, 
and curriculum analytics. These technological developments 
promise to enhance scalability, ensure consistency in 
grading, and provide timely feedback—particularly valuable in 
contexts where human resources are limited.

Despite these advantages, several critical challenges accompany 
AI adoption in education. Key concerns include cultural and 
linguistic bias, fairness, ethical transparency, and the lack 
of robust empirical evidence validating actual educational 
impact. Many AI-driven assessment tools remain test-centric, 
often neglecting broader educational goals such as socio-emotional 
development, creativity, and critical reasoning. This 
limitation is especially pronounced in underrepresented 
languages and cultural contexts.

For the Thai educational landscape, these concerns are even 
more pronounced. Compared to high-resource languages like 
English or Chinese, Thai suffers from limited data 
availability, linguistic diversity across regions, and 
nuanced expressions that challenge direct model generalization. 
As a result, LLM-based systems trained primarily on non-Thai 
data risk perpetuating biases and inconsistencies in grading. 
Such biases can manifest not only in content understanding 
but also in scoring fairness, potentially disadvantaging 
students who express themselves in culturally or linguistically 
distinct ways.

The present thesis directly addresses this gap by developing 
and evaluating an AI-assisted grading framework tailored for 
Thai-language written examinations. By combining LLM-based 
scoring with rule-based reasoning and human-designed rubrics, 
the framework aims to achieve both consistency and fairness. 
Human teachers play a crucial role in refining rubrics and 
validating model performance, ensuring that the automated 
system remains transparent and pedagogically aligned. In 
doing so, this research contributes to the broader discourse 
on equitable AI in education—demonstrating that effective 
integration of LLMs in assessment requires not only technical 
sophistication but also cultural and linguistic sensitivity.

% \cite{Lan2024} discusses the dual role of LLMs in education: 
% on one hand, they support personalized learning, offering 
% tailored feedback and explanations to students; on the other 
% hand, they may encourage academic dishonesty by enabling 
% students to outsource thinking tasks to AI.

\subsection{Teachers’ agency in the era of LLM and generative AI: Designing pedagogical AI agents}

Building upon this perspective, \cite{Lan2024} highlights the 
dual nature of LLMs in educational contexts—acting both as a 
transformative tool for personalized learning and as a 
potential source of academic integrity challenges. On the 
positive side, LLMs can deliver tailored explanations, 
adaptive feedback, and individualized support aligned with 
each learner’s pace and style. However, the same capabilities 
that enable personalization can also encourage academic 
dishonesty, as students may become overly reliant on 
AI-generated outputs rather than engaging in independent 
cognitive effort. This dichotomy underscores the necessity 
for thoughtful pedagogical design and ethical frameworks 
governing how LLMs are integrated into learning environments.

In particular, studies examining teacher and student perceptions 
of generative AI reveal key challenges. Educators report 
difficulty distinguishing between student- and AI-generated 
work, uncertainty about AI’s pedagogical value, and concerns 
that excessive dependence could undermine authentic learning. 
Students, meanwhile, often adopt AI tools superficially, using 
them primarily as search engines or text generators rather 
than as interactive learning companions that foster reflective 
thinking, critical reasoning, or metacognitive development.

To address these issues, the concept of Pedagogical AI Agents 
has been proposed—AI-driven teaching assistants that collaborate 
with human teachers within a team-teaching model. In this 
framework, the teacher remains central, responsible for 
instructional planning, rubric design, and student evaluation, 
while the AI agent delivers adaptive content, real-time 
feedback, and continuous learning analytics. Structured around 
established instructional design models such as 
ADDIE and Dick \& Carey, the approach emphasizes role 
definition, instructional procedures, assessment rubrics, 
guidance and scaffolding, and feedback mechanisms. For example, 
AI agents have been employed to teach sequencing words in 
English, providing structured explanations, contextual examples, 
and interactive exercises with immediate feedback, leading to 
measurable improvements in understanding and retention.

Ultimately, this line of work reinforces that AI should 
complement rather than replace teachers, acting as a collaborative 
co-teacher that extends educators’ reach and efficiency while 
preserving fairness and pedagogical oversight. Within the 
context of the present thesis, this perspective supports 
positioning AI as a partner in Thai-language written exam 
evaluation—ensuring that automation enhances scoring 
consistency and efficiency without compromising human judgment, 
contextual understanding, or ethical responsibility.

% Similarly,\cite{Ryznar2023} warns of potential misuse in assessment contexts, 
% recommending redesigns of exam formats and the use of 
% time-limited conditions to prevent LLM-driven plagiarism or 
% unfair advantages.

\subsection{Exam in the time of chatgpt}

As discussed in \cite{Ryznar2023}, recent research examines 
the implications of LLMs—specifically ChatGPT—on academic 
integrity and assessment design in higher education. The 
extraordinary capabilities of ChatGPT—such as its strong 
performance in complex professional or legal 
examinations—raise fundamental questions about the validity 
of conventional assessment methods. As LLMs become 
increasingly adept at generating coherent, domain-specific 
responses, educators face the risk that student submissions 
may no longer reflect genuine individual understanding or 
reasoning ability.

Several critical concerns are identified. Students may over-rely 
on AI tools to complete assignments or exam tasks, undermining 
independent thinking. Teachers often struggle to distinguish 
AI-generated from human-authored responses, as existing 
plagiarism detection tools are limited. Additionally, online 
examination systems often lack adequate safeguards, 
particularly in remote contexts, threatening fairness and 
reliability in academic evaluation.

To mitigate these risks, multiple strategies are suggested, 
including video-based proctoring, lockdown browsers, 
AI-specific plagiarism detection, and strict time limits. 
However, technical controls alone are insufficient. Effective 
exam design—emphasizing problem-solving, issue-spotting, 
logical reasoning, and analytical synthesis—is recommended as 
the most sustainable and pedagogically sound approach. 
Illustrative examples from legal education show that while 
ChatGPT performs well on definition-based questions, it 
struggles with analytical or case-based scenarios, 
highlighting the importance of measuring higher-order 
cognitive skills rather than mere factual recall.

The study concludes that eliminating AI from education is 
neither feasible nor desirable. Instead, institutions must 
adapt assessment practices to coexist with AI responsibly. A 
multi-layered approach that combines sound exam design, 
fairness-oriented rubrics, and procedural safeguards can 
preserve both integrity and authenticity.

This perspective aligns directly with the present thesis, 
which emphasizes structured rubric-based evaluation and human 
oversight as mechanisms to ensure fairness and validity in 
AI-assisted grading. By incorporating clear teacher-designed 
rubrics, cross-verification between AI and human assessors, 
and transparent reasoning processes, the proposed framework 
addresses the challenges posed by generative AI. In the Thai 
educational context—where linguistic diversity and limited AI 
training data further complicate fairness—well-designed 
assessment structures become particularly essential, 
reinforcing the thesis’s central contribution: an LLM-based 
framework that is pedagogically sound, culturally fair, and 
empirically reliable.

% These studies collectively highlight both the promise and 
% pitfalls of adopting LLMs in educational systems. While they 
% show potential to alleviate teacher workload and improve 
% grading efficiency, they also underscore the importance of 
% trustworthy, transparent, and context-sensitive systems.

\subsection{Summary}

Taken together, these studies illuminate both the transformative 
potential and the inherent complexities of integrating LLMs 
into educational assessment. Collectively, they demonstrate 
that while AI technologies can enhance scalability, efficiency, 
and consistency in grading, their deployment must be guided by 
principles of fairness, transparency, and pedagogical 
integrity. \cite{Holmes2022} establish the foundational 
perspective that AI can democratize access to feedback and 
streamline educational processes, yet they caution that issues 
of cultural and linguistic bias remain unresolved—particularly 
in underrepresented languages such as Thai. \cite{Lan2024} 
further extend this discussion by emphasizing the need for 
human–AI collaboration in the form of pedagogical AI agents, 
which complement rather than replace teachers. Their 
team-teaching model reinforces the idea that successful 
educational AI systems depend on clearly defined human roles, 
robust rubrics, and ethical instructional design.

In contrast, \cite{Ryznar2023} highlights the other side of AI 
integration—the risks associated with academic dishonesty, 
overreliance, and declining authenticity in assessment. His 
findings reveal that without thoughtful redesign, conventional 
exam formats are vulnerable to misuse and may fail to capture 
genuine student learning. Consequently, future assessment 
practices must evolve beyond mechanical grading to emphasize 
reasoning, creativity, and critical thinking—dimensions that 
LLMs can support but not fully replicate.

Across these perspectives, a unifying theme emerges: the 
future of AI-assisted assessment lies in balance. Automation 
should serve as a tool that augments human judgment, not a 
mechanism that replaces it. Effective integration of LLMs 
therefore requires systems that are not only technically 
capable but also contextually grounded—sensitive to linguistic, 
cultural, and pedagogical nuances within each educational environment.

In the context of Thai-language education, where limited data 
resources and unique linguistic structures present additional 
challenges, these insights are particularly pertinent. The 
thesis builds upon this foundation by proposing an AI-assisted 
grading framework that integrates LLM-generated evaluations 
with teacher-designed rubrics and expert verification. This 
approach directly addresses the concerns raised in prior 
research—offering a pathway toward AI assessment that is 
efficient yet fair, data-driven yet human-centered, and 
adaptable to the cultural realities of Thai education. 
These opportunities and challenges are summarized in 
\rfig{opportunities_and_challenges_of_llm}.

\bfigl{.89}{opportunities_and_challenges_of_llm}{Opportunities and Challenges of LLMs in Educational Assessment. Adapted from \cite{Vajjala2025}.}

\section{Leveraging LLMs for Text Understanding and Cognitive Inference}

% While LLMs are often praised for their text generation capabilities, 
% their performance in logical reasoning and semantic inference is less 
% consistent. \cite{Jang2023} demonstrated that GPT-4 performs well in 
% syntactic and linguistic tasks, but struggles with negation, 
% contradiction, and inference-based reasoning. This suggests that while 
% LLMs capture surface-level patterns, deeper cognitive alignment may 
% require additional strategies.

\subsection{Consistency Analysis of ChatGPT}

As shown in \cite{Jang2023}, while LLMs are widely recognized 
for their impressive text generation capabilities, their 
reliability in logical reasoning and semantic inference 
remains inconsistent. Although GPT-4 demonstrates strong 
performance in syntactic and linguistic comprehension tasks, 
it continues to struggle with more complex reasoning forms 
such as negation, contradiction, and inference-based 
decision-making. These findings suggest that while LLMs excel 
at capturing surface-level language patterns, their ability 
to represent deeper cognitive reasoning—particularly those 
involving relational and causal logic—remains limited.

A more systematic investigation of this limitation is 
presented in the Consistency Analysis of ChatGPT (2023), 
which evaluates the trustworthiness and logical consistency 
of ChatGPT and GPT-4 through four categories of reasoning 
tasks: (1) semantic reasoning, which assesses the model’s 
understanding of meaning and contextual relationships; (2) 
negation reasoning, which tests comprehension of oppositional 
statements; (3) symmetric reasoning, where equivalence between 
reversed elements (e.g., “A relates to B” versus “B relates to 
A”) is measured; and (4) transitive reasoning, which examines 
the model’s ability to infer indirect relationships (e.g., if 
A → B and B → C, then A → C).

\begin{figure}[H]
    \centering
    \includegraphics[width=0.49\textwidth]{2.2.1_consistency}
    \caption{Consistency evaluation process of (a) semantic, 
    (b) negation, and (c) symmetric consistency on MRPC and 
    (d) transitive consistency on SNLI, adapted from \cite{Jang2023}.}
    \label{fig:2.2.1_consistency}
\end{figure}

The study revealed that, while GPT-4 and ChatGPT have shown 
measurable improvement in handling negation and transitive 
inference compared to earlier LLM generations, they still fall 
short in maintaining semantic and symmetric consistency. 
Notably, both models exhibited self-contradiction when 
presented with paraphrased inputs or minor variations in 
question order, exposing an inherent input-order sensitivity 
that undermines reliability. These inconsistencies persisted 
even when using larger model architectures, advanced prompting 
strategies, or few-shot learning techniques—suggesting that 
scaling alone does not resolve fundamental reasoning flaws.

Furthermore, the authors caution that continual model enlargement 
is not a sustainable solution, given its escalating economic 
cost and environmental impact. Instead, they argue that the 
future of reliable AI reasoning depends on introducing 
structured, rule-based, and human-guided approaches that 
constrain model outputs within logically coherent boundaries. 
Logical consistency, they contend, is a prerequisite for 
deploying LLMs in high-stakes decision-making domains such as 
law, healthcare, and finance—fields where even subtle 
contradictions can result in serious consequences.

In the context of educational assessment, these insights carry 
critical implications. Logical reliability is central to the 
fairness and transparency of automated grading. If an LLM 
inconsistently interprets similar student responses or fails 
to apply reasoning consistently across equivalent contexts, 
the resulting evaluation cannot be considered trustworthy. 
The present thesis directly addresses this concern by 
integrating rule-based reasoning with LLM evaluation and 
teacher-validated rubrics, ensuring that AI-driven assessments 
remain logically sound and contextually fair. By combining 
structured rubrics and human oversight, the proposed framework 
mitigates the weaknesses identified in existing LLM reasoning 
research—offering a model that aligns machine inference with 
educational judgment in a transparent, reproducible manner.

% \cite{Lin2024} explored model fine-tuning, showing that applying 
% lightweight methods such as LoRA to models like LLaMA-3 allows 
% specialization in tasks such as sentiment intensity regression. 
% Although fine-tuning is outside the scope of this thesis, their 
% findings emphasize the need to align models with specific cognitive 
% tasks rather than expecting universal performance from a single model.

\subsection{Enhanced Sentiment Intensity Regression Through LoRA Fine-Tuning on Llama 3}

As demonstrated in \cite{Lin2024}, beyond reasoning 
consistency, another research direction explores the 
alignment of LLMs with domain-specific cognitive tasks 
through efficient fine-tuning techniques. The study 
investigated Enhanced Sentiment Intensity Regression Through 
LoRA Fine-Tuning on LLaMA-3, aiming to improve the model’s 
sensitivity to emotional gradation rather than mere polarity 
classification. Unlike conventional sentiment analysis, which 
focuses on categorical outputs (positive, negative, or 
neutral), this work targets the continuous prediction of 
emotional strength, providing a more nuanced understanding of 
linguistic affect.

\begin{figure}[H]
    \centering
    \includegraphics[width=0.6\textwidth]{2.2.2_emotion}
    \caption{Comparison of sentiment intensity regression with other 
    sentiment analysis tasks. The intensity regression approach 
    highlights differences between example sentences, whereas 
    classification and emotion detection yield identical categorical 
    outputs. Adapted from \cite{Lin2024}.}
    \label{fig:2.2.2_emotion}
\end{figure}

To achieve this, Lin and colleagues first enhanced the RoBERTa 
model by integrating an Adaptive Weighted Huber Loss and an 
Additive Attention Mechanism, both designed to improve 
prediction precision and contextual focus. The refined outputs 
from this model were then leveraged to fine-tune LLaMA-3 using 
Low-Rank Adaptation (LoRA)—a parameter-efficient technique that 
enables rapid specialization without the computational cost of 
full retraining. Complementary methods, including prompt engineering 
and data augmentation, further strengthened the model’s 
robustness and generalization.

Experimental results on the SemEval 2017 and 2018 benchmarks 
demonstrated that the LoRA-fine-tuned model outperformed 
several baselines—such as XLNet, ALBERT, Gemma-7B, and the 
standard LLaMA-3—in both error reduction and correlation with 
human judgments. Ablation studies confirmed that both the 
custom loss function and the attention mechanism contributed 
critically to these improvements, while text augmentation 
enhanced resilience to linguistic variation. Despite these 
gains, the authors acknowledged the limitations of dataset 
size and computational scale, noting that broader testing with 
larger models and multimodal inputs remains an important 
future direction.

This work underscores a key insight: LLMs cannot excel universally 
across all cognitive dimensions without task-aligned adaptation. 
Rather than treating a model as a one-size-fits-all evaluator, 
fine-tuning allows for the calibration of linguistic sensitivity, 
reasoning focus, and contextual awareness. In the context of 
this thesis, this finding reinforces the rationale for 
designing a structured, rubric-guided, and rule-based 
evaluation framework, where model outputs are aligned with 
teacher-defined cognitive dimensions. Such an approach parallels 
fine-tuning in that it systematically constrains and guides 
the model’s interpretive space—enhancing both the reliability 
and the educational validity of AI-assisted grading.

\begin{figure}[H]
    \centering
    \includegraphics[width=0.8\textwidth]{2.2.3_emotion_arch}
    \caption{Workflow of the Enhanced Sentiment Intensity Regression Model: Input text is processed by the enhanced Roberta model with efficient 
    additive attention and adaptive weighted Huber loss. The output is combined with augmented input data for prompt engineering, followed by 
    fine-tuning of the Llama 3 model using LoRA within the Unsloth framework. Adapted from \cite{Lin2024}.}    
    \label{fig:2.2.3_emotion_arch}
\end{figure}

% In this study, Bloom’s taxonomy is adopted as a classification lens: 
% mapping student responses to different cognitive levels and selecting 
% the most appropriate evaluation strategy (rule-based or LLM-based). 
% This avoids the pitfalls of a one-size-fits-all method and ensures 
% better alignment between assessment goals and AI capabilities.

\subsection{Summary}

Taken together, these studies reveal that while LLMs have made 
substantial progress in linguistic comprehension and contextual 
reasoning, their reliability remains uneven across different 
types of inference and emotional interpretation. Logical 
consistency studies underscore persistent vulnerabilities in 
semantic coherence, symmetry, and contradiction handling, 
while fine-tuning research highlights the potential of 
task-specific adaptation to enhance precision in cognitive and 
affective understanding.

Building on these insights, this thesis adopts Bloom’s taxonomy 
as a guiding cognitive framework to classify and evaluate student 
responses according to their reasoning depth and expressive 
complexity. By mapping answers to distinct cognitive levels, 
the system determines when to employ rule-based reasoning 
(for factual, definitional, or procedural tasks) and when to 
leverage LLM-based evaluation (for interpretive or analytical 
responses). This hybrid design mitigates the “one-size-fits-all” 
limitation of general-purpose models, ensuring that the 
assessment strategy aligns with both the intended learning 
outcomes and the cognitive capabilities of the AI.

Ultimately, this approach bridges the gap between linguistic 
pattern recognition and meaningful reasoning—moving toward 
AI-assisted grading that is not only automated and efficient 
but also cognitively grounded and pedagogically faithful.

\section{LLM Ensembles and Model Selection Strategies}

% Since no single LLM excels universally across all tasks, ensemble 
% strategies have gained attention as a way to mitigate weaknesses. 
% \cite{Sakai2025} introduced QUAD-LLM-MLTC, an ensemble framework 
% for healthcare multi-label classification, which showed significant 
% improvements in semantic accuracy compared to single-model baselines. 

\subsection{QUAD-LLM-MLTC: Large Language Models Ensemble Learning for Healthcare Text Multi-Label Classification}

As demonstrated in \cite{Sakai2025}, no single LLM excels 
universally across all tasks, leading ensemble-based 
strategies to emerge as a promising direction to mitigate 
individual model weaknesses and enhance reliability. The 
study introduced QUAD-LLM-MLTC, an ensemble framework 
designed for healthcare multi-label text classification, 
where semantic ambiguity and data imbalance present 
significant challenges. The framework integrates four 
distinct LLMs—BERT for semantic token extraction, PEGASUS for 
text augmentation, GPT-4o for classification, and BART for 
probabilistic label estimation—within a coordinated pipeline.

By combining outputs through a meta-classifier, the system 
achieved substantial gains in F1-score and robustness compared 
to both traditional machine learning approaches and 
single-model baselines. Notably, the method operated under 
zero-shot conditions, demonstrating that careful orchestration 
of multiple LLMs can yield generalizable and contextually 
aware results even without large annotated datasets. This 
architecture highlights the value of contextual prompt 
engineering, label probability fusion, and multimodel design, 
which collectively enhance semantic precision and consistency 
across tasks.

\begin{figure}[H]
    \centering
    \includegraphics[width=0.89\textwidth]{2.3.1_quad}
    \caption{Overview of the QUAD-LLM-MLTC architecture, illustrating the four-model pipeline for multi-label classification. Adapted from \cite{Sakai2025}.}
    \label{fig:quad}
\end{figure}

The findings from this work underscore a critical insight also 
central to this thesis: reliability and fairness in AI-driven 
assessment depend not solely on model size or capability but 
on structural design and cooperative reasoning. Just as 
QUAD-LLM-MLTC leverages multiple specialized models to balance 
strengths and weaknesses, the proposed AI-assisted grading 
framework in this study similarly integrates rule-based 
reasoning, rubric-guided evaluation, and teacher validation. 
This hybrid structure ensures that automated grading remains 
not only efficient and scalable but also transparent, 
interpretable, and pedagogically aligned with human judgment.

% Similarly, \cite{Tekin2024} proposed LLM-TOPLA, which clusters model 
% outputs to enhance robustness and efficiency in processing.

\subsection{LLM-TOPLA: Efficient LLM Ensemble by Maximising Diversity}

As demonstrated in \cite{Tekin2024}, LLM-TOPLA is a 
diversity-optimized ensemble framework designed to enhance 
both robustness and efficiency across a variety of reasoning 
and generation tasks. Unlike conventional ensemble methods 
that prioritize averaging model predictions, LLM-TOPLA 
explicitly measures and leverages error diversity among base 
models to form more balanced and resilient outcomes. Its 
architecture integrates three core innovations: a Focal 
Diversity Metric for quantifying complementary model 
behavior, a Diversity-Optimized Ensemble Pruning Algorithm 
for selecting high-performing sub-ensembles, and a 
Learn-to-Ensemble mechanism that adapts output fusion 
strategies to task type—TOPLA-Weighted for reasoning tasks 
and TOPLA-Summary for generative ones.

\begin{figure}[H]
    \centering
    \includegraphics[width=0.95\textwidth]{2.3.2_llm_topla}
    \caption{Overview of the LLM-TOPLA framework, illustrating diversity-based ensemble selection and adaptive fusion mechanisms. Adapted from \cite{Tekin2024}.}
    \label{fig:2.3.2_llm_topla}
\end{figure}

Evaluated across six benchmark datasets—MMLU, GSM8k, SearchQA, 
XSum, BBH, and ARC—the framework achieved consistent performance 
gains over leading ensemble methods such as Mixtral, More 
Agents, and LLM-Blender. Results showed 2–8\% accuracy improvements 
in reasoning benchmarks and over 30\% increases in F1 and 
ROUGE scores for generative tasks, all while maintaining high 
efficiency with fewer sub-models. This demonstrates that fostering 
model diversity and structured output fusion can outperform even 
large-scale fine-tuning approaches, offering a scalable path 
toward reliable LLM orchestration.

The insights from LLM-TOPLA directly parallel the hybrid 
philosophy of this thesis. Both emphasize that structured 
diversity—whether through model ensembles or rubric-guided 
reasoning—serves as the foundation of fairness and stability 
in AI-driven evaluation. In the proposed grading framework, 
this principle manifests through the integration of rule-based 
reasoning, rubric alignment, and teacher moderation, ensuring 
that no single evaluation component dominates. By balancing 
automated inference with human-informed structure, the system 
aspires to deliver assessments that are robust, interpretable, 
and pedagogically consistent across cognitive levels and 
linguistic variations.

% These works demonstrate that strategic model selection or combination 
% of models can enhance reliability in real-world applications. Inspired 
% by this, the current study adopts a hybrid strategy: selecting the 
% best-performing LLM for each Bloom’s taxonomy category while 
% integrating rule-based reasoning for tasks that lend themselves to 
% deterministic evaluation.

\subsection{Summary}

These studies collectively demonstrate that strategic model 
selection and structured ensemble design can substantially 
improve the reliability, robustness, and interpretability of 
LLM-based systems. By combining models with complementary 
strengths or integrating diversity-aware mechanisms, 
researchers have achieved more stable and context-sensitive 
outcomes across domains.

Building on this insight, the present study employs a hybrid 
evaluation strategy: selecting the best-performing LLM for 
each Bloom’s taxonomy category while incorporating rule-based 
reasoning for tasks that require deterministic and explainable 
assessment. This approach balances the adaptability of language 
models with the rigor of rule-based logic, ensuring evaluations 
that are not only accurate and fair, but also transparent and 
pedagogically grounded—a crucial requirement for Thai-language 
educational contexts.

%%%%%%%%%%%%%%%%%%%%%%%%%%%%%%%%%%%%%%%%%%%%%%%%%%%%%%%%%%%

\section{Bloom’s Taxonomy in Educational AI and Cognitive Classification}

Bloom’s taxonomy provides a structured theoretical foundation for 
aligning assessment with cognitive skill levels. It divides learning 
objectives into six hierarchical categories, each representing a 
progressive stage of cognitive development. The taxonomy not only 
guides curriculum design and learning objectives but also serves as a 
framework for evaluating the depth of understanding demonstrated in 
student responses—making it particularly useful for AI-supported 
assessment systems.

\begin{enumerate}
    \item \textbf{Remembering} – recalling factual knowledge and basic concepts.  
    This foundational level involves recognizing, identifying, and retrieving information from memory. Assessments targeting this level typically include direct recall or multiple-choice questions that measure knowledge retention rather than comprehension.

    \item \textbf{Understanding} – explaining ideas or concepts.  
    Learners at this level demonstrate comprehension by interpreting, summarizing, or classifying information in their own words. Evaluation focuses on whether students can grasp meaning and translate it across contexts rather than merely restating facts.

    \item \textbf{Applying} – using knowledge in new contexts.  
    This level reflects the ability to transfer theoretical knowledge to practical or situational problems. Students apply learned concepts to execute procedures, solve problems, or demonstrate techniques, making it a key stage for performance-based assessment.

    \item \textbf{Analyzing} – breaking down information into components and recognizing relationships or patterns.  
    Analytical thinking involves distinguishing between evidence and argument, identifying cause-and-effect relationships, and examining underlying assumptions. This level assesses critical reasoning and logical structure in student responses.

    \item \textbf{Evaluating} – making judgments or justifying decisions based on criteria or standards.  
    Learners are expected to critique, defend, or assess ideas using well-defined criteria. At this stage, assessment measures reasoning quality, justification of opinions, and the ability to evaluate alternatives—key for rubric-based evaluation systems.

    \item \textbf{Creating} – producing new or original work by integrating knowledge and creativity.  
    The highest cognitive level involves synthesizing ideas to form new solutions, models, or perspectives. This level tests originality, innovation, and the ability to generate coherent, meaningful output—making it a strong indicator of higher-order thinking.
\end{enumerate}

Each level in Bloom’s taxonomy reflects increasing cognitive complexity, from surface-level recall to deep conceptual reasoning and creativity. 
When applied to automated grading, these categories allow AI systems to align evaluation criteria with learning objectives, ensuring that 
the assessment process captures both the \emph{depth} and \emph{quality} of student understanding. 
This hierarchical structure is especially valuable in this thesis, serving as the cognitive backbone for mapping responses to appropriate evaluation strategies—rule-based for lower levels and LLM-assisted for higher-order reasoning.

% This taxonomy has been widely adopted in educational AI research. 
% \cite{Abduljabbar2015} applied machine learning to classify 
% questions into Bloom’s levels, showing its utility for structuring 
% automated evaluation. 

\subsection{Exam questions classification based on Bloom’s taxonomy cognitive level using classifiers combination}

As demonstrated in \cite{Abduljabbar2015}, Bloom’s taxonomy 
has been widely adopted in educational AI research as a 
framework for classifying both questions and student 
responses according to cognitive complexity. Their work 
showed the effectiveness of using machine learning to 
automatically categorize questions into Bloom’s levels, 
highlighting its value in structuring automated evaluation 
systems.

{
\setlength{\LTleft}{-1.7cm}
\setlength{\LTright}{0pt}
\begin{longtable}{l p{0.3\textwidth} p{0.3\textwidth} p{0.3\textwidth}}
    \caption{Six Categories of the Cognitive Domain Based on the Revised Bloom’s Taxonomy (2001) With Illustrative Examples}
    \label{tab:bloom}
    \\

    % For Header
    \hline
    \textbf{Category} & \textbf{Sample Keywords} & \textbf{Sample Behaviours} & \textbf{Sample Questions} \\
    \hline
    \endfirsthead

    % For Overflow Header
    \hline
    \textbf{Category} & \textbf{Sample Keywords} & \textbf{Sample Behaviours} & \textbf{Sample Questions} \\
    \hline
    \endhead

    % For Footer
    \hline
    \endfoot

    Remembering & Define, list, name, recall, recognize & Students are able to define the 6 levels of Bloom's taxonomy & Recall basic facts and terms. \\
    \cline{1-4}
    Understanding & Classify, explain, describe, discuss, explain, summarize. & Students can explain the purpose of Bloom's taxonomy. & Explain relationships among concepts. \\
    \cline{1-4}
    Applying & Demonstrate, implement, solve, use, execute & Students write an instructional objective for a method in the program. & Apply knowledge to new situations \\
    \cline{1-4}
    Analyzing & Analyze, compare, contrast, differentiate, examine. & Students compare and contrast cognitive and affective domains. & Break concepts into parts and examine relationships. \\
    \cline{1-4}
    Evaluating & Assess, critique, defend, judge, validate. & Students judge the effectiveness of Revised Bloom's taxonomy. & Justify the concept of inheritance and give code examples. \\
    \cline{1-4}
    Creating & Assemble, construct, design, formulate, invent. & Students design a classification method for educational objectives. &  Combine elements to form a coherent whole or original product. \\
    
\end{longtable}
}

Recent work on automatic exam question classification further 
advances this line of research by addressing the inherent 
difficulty of mapping short and sparse textual data to 
cognitive levels. Because exam questions are typically 
concise and context-limited, traditional classification 
methods often suffer from low accuracy and poor generalization.

To overcome these challenges, researchers have proposed a 
Voting Ensemble Model that combines the strengths of multiple 
classifiers—namely Support Vector Machine (SVM), Naïve Bayes 
(NB), and k-Nearest Neighbour (k-NN). In addition, feature 
selection techniques such as Chi-Square, Mutual Information, 
and Odds Ratio were employed to isolate the most informative 
features. Experimental results showed that the ensemble 
approach outperformed all individual models, with the best 
accuracy achieved using Mutual Information and 250 selected 
features.

\begin{figure}[H]
    \centering
    \includegraphics[width=0.7\textwidth]{2.4.1_system}
    \caption{Overview of the proposed classifier combination model for Bloom’s-level prediction. Adapted from \cite{Abduljabbar2015}.}
    \label{fig:2.4.1_system}
\end{figure}

These findings highlight that ensemble learning enhances the 
accuracy, stability, and scalability of Bloom’s-level 
classification, particularly in handling short-text educational 
data. This approach provides a robust foundation for fine-grained 
exam analysis and adaptive assessment design.

From the perspective of this thesis, such work directly 
supports the integration of Bloom’s taxonomy as the structural 
backbone of the proposed AI-assisted grading framework. It 
also reinforces the rationale for employing machine learning 
and ensemble-based strategies to achieve more precise and 
cognitively aligned evaluation of exam content.

% \cite{Rodrigo2015} emphasized the importance 
% of inter-rater agreement and gold-standard datasets, both critical for 
% training reliable assessment systems.

\subsection{On Evaluating the Contribution of Validation for Question Answering}

\cite{Rodrigo2015} emphasized the importance of inter-rater 
agreement and gold-standard datasets, both critical for 
training reliable assessment systems.
Building on this, subsequent research has examined how validation 
processes influence the development and reliability of Question 
Answering (QA) systems, especially in the context of answer 
validation (AV) modules that verify the correctness of responses.

Traditional evaluation metrics such as Accuracy, Precision, 
Recall, and F1-score often fail to reflect the real-world 
effectiveness of AV modules. These intrinsic measures assess 
a module’s internal consistency but do not indicate whether 
it actually improves the overall performance of the QA system.

To address this limitation, researchers proposed a QA-based 
evaluation framework that assesses AV modules based on their 
extrinsic impact—that is, their contribution to the end-to-end 
performance of the full QA system. The framework introduces 
task-aligned metrics such as Mean Reciprocal Cost (MRC) and 
Macro-Precision, evaluated using the AVE 2008 dataset across 
both simulated and real QA systems.

Experimental results revealed that conventional evaluation 
methods could produce misleading conclusions in up to 88\% of 
cases, whereas the proposed QA-based framework provided a more 
accurate reflection of true system performance. This evidence 
strongly supports the view that AV modules should not be 
assessed in isolation but as integrated components within the 
broader QA pipeline.

The study therefore argues for a shift from intrinsic to 
extrinsic evaluation, emphasizing that effective system 
validation must consider context, fairness, and real-world 
impact rather than raw numerical performance alone.

This perspective aligns closely with the approach adopted in 
this thesis. The proposed AI-assisted grading framework emphasizes 
context-aware evaluation that reflects genuine assessment 
outcomes and fairness. By incorporating structured rubrics, 
rule-based reasoning, and teacher validation, the system 
ensures that LLM evaluations of Thai-language exam responses 
are not only accurate but also transparent, interpretable, and 
pedagogically meaningful.

% More recent studies explore hybrid approaches. \cite{Dodia2023} 
% combined rule-based methods with machine learning to grade short-answer 
% responses, 

\subsection{Machine Learning-based Automated System for Subjective Answer Evaluation}

As demonstrated in \cite{Dodia2023}, more recent studies 
explore hybrid approaches that combine rule-based methods 
with machine learning techniques to grade short-answer 
responses, aiming to reduce the time, cost, and workload 
associated with manual grading.

The proposed framework integrates three key components:

\begin{enumerate}
    \item BERT embeddings – to capture the semantic meaning of both teacher and student answers;
    \item Cosine similarity – to measure semantic closeness between vectorized responses; and
    \item Jaccard distance – to assess lexical overlap between word sets in answers.
\end{enumerate}

The similarity scores from these three techniques were weighted 
(3.5 : 1 : 0.5) and passed into a neural network with a sigmoid 
activation function to produce the final student score.

\begin{figure}[H]
    \centering
    \includegraphics[width=0.8\textwidth]{2.4.3_ml_quiz}
    \caption{Schematic architecture of the proposed automated grading system, integrating BERT embeddings, cosine similarity, and Jaccard distance. Adapted from \cite{Dodia2023}.}    
    \label{fig:2.4.3_ml_quiz}
\end{figure}

The dataset consisted of 100 exam questions, each paired with 
a model answer provided by a teacher and five student responses 
per question. Experimental results demonstrated that combining 
multiple similarity metrics achieved the best performance, with 
an accuracy of 91\%, precision of 83\%, recall of 79\%, and 
an error rate of 9.01\%. Compared to baseline methods such as 
LSTM, Naïve Bayes, and KNN, this approach produced scores that 
more closely aligned with teacher evaluations.

Overall, the study showed that combining semantic and statistical 
similarity measures within a neural network framework can 
significantly improve the reliability and accuracy of automated 
grading for descriptive answers.

This work directly supports the direction of the present 
thesis, illustrating that integrating semantic models with 
rule-based or similarity-based scoring methods enhances 
grading accuracy, fairness, and interpretability—key principles 
underpinning the design of the AI-assisted Thai exam evaluation 
framework.

% while \cite{Hubbard2017} demonstrated that question formats 
% (e.g., multiple-true-false vs. free-response) influence the depth of 
% student reasoning captured.


\subsection{How Question Types Reveal Student Thinking: An Experimental Comparison of Multiple-True-False and Free-Response Formats}

While \cite{Hubbard2017} demonstrated that question formats 
(e.g., multiple-true-false vs. free-response) influence the 
depth of student reasoning captured. This study investigated 
how different exam formats—particularly 
Multiple-True-False (MTF) and Free-Response (FR)—affect 
students’ conceptual understanding and reasoning processes 
in undergraduate biology courses.

Using a crossover experimental design with over 400 students 
(\rfig{fig:crossover}), the researchers compared 
performance, conceptual understanding, 
and diagnostic value across the two question types. Results 
showed a correlation in overall performance between MTF and FR 
formats, though students generally scored higher on MTF items.

\begin{figure}[H]
    \centering
    \includegraphics[width=0.7\textwidth]{2.4.4_crossover}
    \caption{Crossover experimental design comparing MTF and 
    FR question formats in undergraduate biology exams. 
    Adapted from \cite{Hubbard2017}.}
    \label{fig:crossover}
\end{figure}

MTF questions revealed mixed conceptions—students expressing 
both correct and incorrect ideas simultaneously—while FR responses 
contained fewer misconceptions but tended to be ambiguous or 
vague, particularly among lower-performing students.

In terms of scoring, partial-credit MTF formats tended to 
inflate scores, while all-or-nothing MTF yielded results more 
consistent with FR scoring.

Overall, the study found that MTF questions are effective for 
quickly identifying specific misconceptions, whereas FR 
questions better capture deeper cognitive reasoning, albeit 
with greater scoring complexity and subjectivity. The authors 
concluded that hybrid exam formats can provide a practical 
balance between diagnostic precision, grading efficiency, and 
authentic reflection of student understanding.

This finding aligns closely with the present thesis, emphasizing 
that exam design directly shapes assessment validity and 
interpretability. In the context of AI-assisted grading, this 
underscores the need for structured rubrics and teacher 
validation to ensure that automated scoring of open-ended 
responses maintains both efficiency and fairness.
Conceptually, MTF mirrors AI-based scoring—fast and 
consistent yet prone to bias, while FR parallels human 
grading—deeper but subjective. The proposed hybrid framework 
in this thesis—LLM-based scoring combined with rubric 
reasoning and teacher validation—therefore represents an 
ideal synthesis of both approaches, achieving balanced, 
interpretable, and trustworthy evaluation outcomes.

% These findings support the argument that Bloom’s taxonomy is not only 
% a theoretical tool but also a practical framework for guiding 
% AI-supported assessment. By aligning AI evaluation methods with 
% cognitive levels, systems can achieve greater accuracy, transparency, 
% and interpretability.

\subsection{Summary}

Collectively, these studies demonstrate that Bloom’s taxonomy 
provides not only a theoretical scaffold for understanding 
cognitive complexity but also a practical foundation for 
designing and validating AI-assisted assessment systems. By 
classifying exam questions and student responses along 
cognitive dimensions, researchers have shown improvements in 
grading accuracy, fairness, and interpretability. Ensemble and 
hybrid machine learning models further enhance this process, 
effectively addressing challenges in short or ambiguous responses 
through semantic and rule-based integration. Moreover, studies 
on question formats highlight that assessment design itself 
influences the depth and validity of reasoning captured, 
underscoring the importance of aligning question structure 
with evaluation strategy.

Building on these insights, the present study applies Bloom’s 
taxonomy as a classification and decision framework—mapping 
responses to appropriate cognitive levels and selecting the 
most suitable evaluation method, whether LLM-based, rule-based, 
or teacher-validated. This alignment ensures that AI-assisted 
grading achieves not only computational efficiency but also 
transparency, pedagogical relevance, and cognitive fidelity—key 
requirements for trustworthy educational assessment in Thai-language 
contexts.


% \begin{center}
% \begin{longtable}{l p{4cm} p{7cm}}    % in case want separate line \begin{longtable}{|l|p{4cm}|p{7cm}|}

% \caption{Six Categories of the Cognitive Domain Based on the Revised Bloom’s Taxonomy (2001) With Illustrative Examples}
% \label{tab:revised-bloom} 
% \\

% % For Header
% \hline
% \textbf{Category} & \textbf{Sample Keywords} & \textbf{Sample Behaviours / Sample Questions} \\
% \hline
% \endfirsthead

% % For Overflow Header
% \hline
% \textbf{Category} & \textbf{Sample Keywords} & \textbf{Sample Behaviours / Sample Questions} \\
% \hline
% \endhead

% % For Footer
% \hline
% \endfoot

% Remembering &
% Define, list, name, recall, recognize &
% Students are able to define the 6 levels of Revised Bloom’s Taxonomy. \\
% \cline{2-3}         % Draws a horizontal line only under columns 2 and 3
%  & \textbf{Knowledge:} & Recall basic facts and terms. \\
% \cline{2-3}         % Draws a horizontal line only under columns 2 and 3
%  & \textbf{Sample Questions:} & Define each level of Revised Bloom’s Taxonomy. \\
% \hline

% Understanding &
% Classify, describe, discuss, explain, summarize &
% Students can explain the purpose of the Revised Bloom’s Taxonomy. \\
% \cline{2-3}
%  & \textbf{Knowledge:} & Explain relationships among concepts. \\
% \cline{2-3}
%  & \textbf{Sample Questions:} & Explain the purpose and importance of each level. \\
% \hline

% Applying &
% Demonstrate, implement, solve, use, execute &
% Students write an instructional objective for a method in the program. \\
% \cline{2-3}
%  & \textbf{Knowledge:} & Apply knowledge to new situations. \\
% \cline{2-3}     
%  & \textbf{Sample Questions:} & Demonstrate the structure of each level in a program. \\
% \hline

% Analyzing &
% Analyze, compare, contrast, differentiate, examine &
% Students compare and contrast cognitive and affective domains. \\
% \cline{2-3}
%  & \textbf{Knowledge:} & Break concepts into parts and examine relationships. \\
% \cline{2-3}
%  & \textbf{Sample Questions:} & List advantages/disadvantages of using \texttt{ArrayList} vs array. \\
% \hline

% Evaluating &
% Assess, critique, defend, judge, validate &
% Students judge the effectiveness of Revised Bloom’s Taxonomy. \\
% \cline{2-3}
%  & \textbf{Knowledge:} & Make judgments based on criteria and standards. \\
% \cline{2-3}
%  & \textbf{Sample Questions:} & Justify the concept of inheritance and give code examples. \\
% \hline

% Creating &
% Assemble, construct, design, formulate, invent &
% Students design a classification method for educational objectives. \\
% \cline{2-3}
%  & \textbf{Knowledge:} & Combine elements to form a coherent whole or original product. \\
% \cline{2-3}
%  & \textbf{Sample Questions:} & Create a new method integrating cognitive, affective, and psychomotor domains. \\
% \hline

% \end{longtable}
% \end{center}

% \section{Summary of Literature Review}

% The reviewed literature highlights three key insights relevant to 
% this thesis:

% \begin{itemize}
%     \item \textbf{LLMs in Education}: They hold significant potential for automating grading and feedback, but face challenges of fairness, cultural bias, and academic integrity.
%     \item \textbf{Text Understanding and Reasoning}: LLMs demonstrate strengths in syntax and semantics but remain inconsistent in logical reasoning, supporting the need for hybrid frameworks.
%     \item \textbf{Ensemble and Cognitive Alignment}: Combining multiple LLMs or matching models to cognitive levels, guided by Bloom’s taxonomy, enhances reliability and interpretability.
% \end{itemize}

% Building upon these insights, this thesis proposes a hybrid AI-assisted 
% grading framework tailored to Thai written exams. The framework 
% integrates LLMs with rule-based reasoning, structured under Bloom’s 
% taxonomy, to ensure fair, consistent, and pedagogically aligned 
% evaluation.

\section{Summary of Literature Review}

The reviewed literature collectively reveals four major insights that form the conceptual foundation of this thesis:

\begin{itemize}
    \item \textbf{AI and LLMs in Education}: Large Language Models (LLMs) have transformed educational technology by enabling scalable, personalized feedback and automated grading. However, scholars emphasize persistent challenges related to fairness, transparency, linguistic bias, and academic integrity, underscoring the need for teacher-centered and ethically guided integration.
    \item \textbf{Assessment Design and Validation}: Studies highlight that evaluation systems must go beyond technical accuracy to measure real-world effectiveness. Reliable assessment requires validation through inter-rater agreement, rubric alignment, and extrinsic performance metrics—ensuring that automated grading truly reflects human educational values.
    \item \textbf{Automated Question and Response Classification}: Research on question classification and short-answer grading demonstrates that hybrid and ensemble learning models outperform single-model or purely rule-based systems. These approaches handle short and sparse educational texts more effectively, improving both precision and adaptability in cognitive-level prediction.
    \item \textbf{Cognitive Frameworks and Bloom’s Taxonomy Alignment}: Bloom’s taxonomy provides a crucial theoretical lens for aligning AI-driven assessment with cognitive skill levels. Empirical findings suggest that linking evaluation criteria to these hierarchical levels enhances interpretability, fairness, and pedagogical validity.
\end{itemize}

Synthesizing these perspectives, this thesis proposes a hybrid 
AI-assisted grading framework tailored for Thai-language written 
examinations. The framework integrates LLM-based evaluation with 
rule-based reasoning and teacher-verified rubrics, all structured 
under Bloom’s taxonomy. This design aims to achieve accurate, 
transparent, and culturally fair assessment—bridging the gap 
between AI efficiency and educational authenticity.
